\scanpage{11}
5. Relation avec les ``noyaux'' de Mac-Lane, ou ``liens''. Soit $\unsure{G|C}$ groupe, $F$ groupe,
\[
\varphi\colon G\longrightarrow\operatorname{Aut\,ext}(F)
\]
hom (\unsure{rev. à|revient à|cela $=$} se donner un lien sur $B_G$). On va lui associer une \inserted{Gr-}catégorie \deleted{de Picard} $\unsure{\mathscr P|P}(\varphi)$ d'invariants $G$ et $\deleted{\unsure{M}}\ M=\mathrm Z(F)$ (centre)

a) \uline{Description directe}. On a la \inserted{Gr} catégorie $\mathscr P(F)$ des bitorseurs sous $F$, qui a comme invariants
\[
\begin{aligned}
\pi_0(\mathscr P(F))&\in\operatorname{Aut\,ext}F\\
\pi_1(\mathscr P(F))&=\mathrm Z(F)\qquad\text{centre de }F
\end{aligned}
\]
On définit donc $\mathscr P(\varphi)$ comme image inverse.

\deleted{b) Description en termes des 2-gerbes liées par $M$.}

\uline{NB} Les trivialisations de $\mathscr P(\varphi)$ correspondent aux réalisations de $\varphi$ par une extension de $G$ par $F$ subordonnée à $\varphi$. Il faudrait encore préciser ce point en termes \unsure{de|des} l'équivalence de 2-catégories (2-torseurs sous une 2-catégorie de Picard). On \unsure{retrouve|retrouvera} ainsi la classe de Mac-Lane dans $H^3(G,M)$, obstruction à la réalisabilité du ``noyau'' $\varphi$ par une extension (ou encore du lien $\varphi$ par une 1-gerbe). \deleted{\unsure{Notons}}\inserted{Mais} \deleted{\unsure{remarquons}}\inserted{le fait que} \deleted{\unsure{tous}} \unsure{élément|élém\textsuperscript{t}|éléments} de $H^3(\unsure{C|G},M)$ (i.e. tte catégorie de Picard \unsure{épinglée|épinglées} par $(G,M)$) puisse ainsi s'obtenir, à l'aide d'un ``noyau'', ne sort pas \inserted{(pour l'instant)} \unsure{d'un fait général|d'un résultat général|d'une logique générale}.
% archive sheet number: 10
